\section{Problem statement and the inclusion function}
\label{sec:prelim}

The collision queries treated here reduce to finding roots of a map $\Fdiff$ on
a three-dimensional parameter box, and that map is multilinear. Multilinearity
means that the eight corner values of a box give the \emph{exact} range of
$\Fdiff$ over it, which in real arithmetic would settle the matter. In floating
point the corner values are perturbed, and a certified bound on the perturbation
is what stands between the search and an unsound answer. This section sets out
both properties, together with the tolerances that tell the search when to stop.
\Cref{sec:guarantee} turns them into a guarantee, and \cref{sec:parallel} into
two implementations.

\subsection{The queries and their difference function}
\label{sec:prelim:queries}

A mesh is given at two instants. Vertex $p$ is at $p^{0}$ at time $t=0$ and at
$p^{1}$ at $t=1$, and moves affinely in between, $p(t) = (1-t)\,p^{0} +
t\,p^{1}$. Two triangle meshes intersect during the step if and only if a vertex
passes through a face or an edge crosses an edge~\citep{provot1997collision}, so
the whole problem reduces to two elementary queries, each expressed as a root of
a map $\Fdiff : [0,1]^{3} \to \Rthree$.

For a \emph{vertex-face} query with vertex $\primA{q}$ and face
$(\primB{f_{1}},\primB{f_{2}},\primB{f_{3}})$, writing $o = 1-u-v$ for the third
barycentric coordinate,
\begin{equation}
\label{eq:fvf}
\Fdiff(t,u,v) \;=\; \primA{q(t)} \;-\; \primB{\bigl[\,o\,f_{1}(t) + u\,f_{2}(t)
+ v\,f_{3}(t)\,\bigr]}, \qquad (u,v)\in\Delta,
\end{equation}
where $\Delta = \{(u,v) : u,v \ge 0,\ u+v \le 1\}$. For an \emph{edge-edge}
query with edges $(\primA{a_{0}},\primA{a_{1}})$ and
$(\primB{b_{0}},\primB{b_{1}})$,
\begin{equation}
\label{eq:fee}
\Fdiff(t,u,v) \;=\; \primA{\bigl[a_{0}(t) + u\,(a_{1}(t)-a_{0}(t))\bigr]} \;-\;
\primB{\bigl[b_{0}(t) + v\,(b_{1}(t)-b_{0}(t))\bigr]}, \qquad (u,v)\in[0,1]^{2}.
\end{equation}
A contact occurs exactly when $\Fdiff$ vanishes somewhere on the domain, and the
earliest time of impact is $\toitrue = \min\{\,t : \Fdiff(t,u,v)=0 \text{ for
some admissible } (u,v)\,\}$.

Quad meshes add a third query of the same shape. For a vertex against a
bilinear quad with corner nodes $\primB{g_{1}},\dots,\primB{g_{4}}$,
\begin{equation}
\label{eq:fvq}
\Fdiff(t,u,v) \;=\; \primA{q(t)} \;-\; \primB{\sum_{k=1}^{4} w_{k}(u,v)\,
g_{k}(t)}, \qquad
w = \bigl((1-u)(1-v),\; u(1-v),\; (1-u)v,\; uv\bigr),
\end{equation}
on $(u,v)\in[0,1]^{2}$. The weights in \cref{eq:fvq} put the nodes in
\emph{lexicographic} order, the corners $(0,0)$, $(1,0)$, $(0,1)$ and $(1,1)$ of
the parameter square, and not the cyclic winding a quad is usually stored with.
\Cref{sec:np:quad} quantifies the cost of an incorrect ordering.

The three queries share the property that the search exploits. Each of
\cref{eq:fvf,eq:fee,eq:fvq} is \emph{multilinear}: affine in $t$ for fixed
$(u,v)$, affine in $u$ for fixed $(t,v)$, and affine in $v$ for fixed $(t,u)$.
On an axis-aligned box $\dom =
[t_{\ell},t_{h}]\times[u_{\ell},u_{h}]\times[v_{\ell},v_{h}]$ a multilinear map
attains its extrema at vertices of the box, so for each component $d$,
\begin{equation}
\label{eq:hull}
\Fdiff_{d}(\dom) \;=\; \bigl[\;\min_{c \in V(\dom)} \Fdiff_{d}(c),\;
\max_{c \in V(\dom)} \Fdiff_{d}(c)\;\bigr],
\end{equation}
where $V(\dom)$ are the eight corners. \Cref{eq:hull} is an equality and not an
enclosure: evaluating the eight corners and taking the componentwise hull gives
the exact range of $\Fdiff$ over $\dom$. That property distinguishes the
inclusion function of \citet{wang2021tight} from generic interval arithmetic,
which widens at every operation, and it allows a box to be discarded on the
strength of its corners alone.

\subsection{Error bound and the tolerances}
\label{sec:prelim:err}

The exactness of \cref{eq:hull} is a statement about real arithmetic, whereas
the eight values a machine computes are $\Fdiff_{d}(c)$ perturbed by rounding.
Since \cref{eq:hull} is the only thing licensing a box to be discarded, the size
of that perturbation is the only thing standing between the search and a lost
root.

\citet{wang2021tight} bound it by a forward error analysis in the standard
style~\citep{higham2002accuracy}, the same reasoning that underlies the
floating-point filters used for exact geometric
predicates~\citep{attene2020indirect}. With $\mach$ the unit roundoff and $M$
the largest absolute coordinate appearing among the eight input points of the
query,
\begin{equation}
\label{eq:err}
\delta \;=\; \min(M,\,1),
\qquad
\epsilon_{d} \;=\; \kappa\,\mach\,\delta^{3},
\qquad
\kappa = \begin{cases}
30 & \text{vertex-face},\\
28 & \text{edge-edge},
\end{cases}
\end{equation}
and every computed corner value lies within $\epsilon_{d}$ of the exact one.

Two details of \cref{eq:err} are essential. The first is that $\delta$ is $M$
clamped from \emph{above} by $1$. Clamping from below instead would make the pad
grow as the cube of the scene's diameter, reaching $10^{6}$ at a scale of $100$,
which the acceptance test would then have to spend.

The second is that the bound is derived for one particular sequence of
floating-point operations. An algebraically identical but differently associated
expression has a different rounding profile, and $\kappa$ does not cover it. The
corner evaluation is therefore written in TightInclusion's order and is not free
to be rearranged. The implementation pins IEEE semantics across those regions,
and no \texttt{-ffast-math} compilation flag is used anywhere in this work.

No constant is certified in the literature for the bilinear query of
\cref{eq:fvq}. We use $\kappa = 38$, which exceeds both certified constants and
is therefore safe in the direction that matters, and we record it as an
over-estimate; lowering it requires a derivation.

The bound determines when a box may be discarded, and a second quantity
determines when a box is small enough for refinement to stop. That quantity is
measured per axis, since the three axes of the domain are not commensurate: a
width in $t$ and a width in $u$ mean different things. Following
\citet{wang2021tight}, a user-facing distance tolerance $\tau$ in the codomain
becomes per-axis domain tolerances by division by the rate at which $\Fdiff$
changes along each axis,
\begin{equation}
\label{eq:tol}
\tau_{0} = \frac{\tau}{L_{t}},\qquad
\tau_{1} = \frac{\tau}{L_{u}},\qquad
\tau_{2} = \frac{\tau}{L_{v}},
\end{equation}
with each $L$ the largest $\ell_{\infty}$ displacement of the relevant corner
differences along that axis.

\Cref{eq:tol} requires a cap. A query whose geometry barely moves along one axis
divides by nearly zero and receives an effectively infinite tolerance, so that
every box satisfies ``width $\le$ tolerance'', the root box is accepted on
sight, and the query reports $\toi = 0$. That outcome is conservative, being
early and not late, and it is also uninformative. The quotients are therefore
clamped at $\tau_{0}\le 10^{-3}$ and $\tau_{1},\tau_{2}\le 10^{-2}$, matching
the reference. The same two caps are compiled into the device kernels, so both
processors stop in the same place.

Together, \cref{eq:hull,eq:err,eq:tol} fix what an implementation of this family
can promise. Exactness is not attainable at simulation speed, whereas
conservativeness in the sense of \cref{def:conservative} is, at the price of
false positives, which \cref{sec:results} reports per scene as the cost they
are.

Deciding a query costs far more than deciding whether it is worth asking, so the
pipeline applies the cheaper test first. \Cref{sec:bp} describes that filter,
which needs none of the machinery above, and \cref{sec:np} the search that uses
all of it.
